\section{Experimental setup}
\label{sec:experiments}

\paragraph{Tool pool.}
We curate 23 MCP servers (151 tool endpoints) spanning seven categories: \texttt{arxiv}, \texttt{excel}, \texttt{fetch}, \texttt{filesystem}, \texttt{pdf}, \texttt{search}, and \texttt{wikipedia}. Servers are smoke-tested for successful initialization and a well-formed response before entering the pool. Four synthetic failure modes (\texttt{truncated}, \texttt{hollow}, \texttt{flaky}, \texttt{stale}) are injected, one variant per category, producing 28 additional endpoints alongside the real tools to give the reranker ground-truth corruptions to learn.

\paragraph{Training queries.}
Queries are generated by an LLM, with $\sim$175 per category across four paraphrase templates. Success labels are produced by running each query through the pipeline and judging the tool output with a calling-model-matched judge (Section~\ref{sec:method}).

\paragraph{Online training.}
Training proceeds over 1{,}200 queries, round-robin across six LLMs (\texttt{gpt-5.4-nano}, \texttt{grok-4.1-fast}, \texttt{gemini-3.1-flash-lite}, \texttt{gemma-4-26b}, \texttt{qwen3.5-flash}, \texttt{deepseek-v3.2}). The reranker retrains every 50 feedback records. A retriever-only bootstrap pass is used to avoid cold-start policy collapse.

\subsection{GAIA ground-truth benchmark}
\label{sec:gaia-gt}

GAIA tasks are multi-hop and underspecified for the tool, so we cannot use them to measure tool routing directly. We convert each GAIA level 1--2 task into a single-hop probe over the same artifact with a verifiable expected substring:

\begin{quote}\small
\textbf{Original (GAIA, level 2):} ``The attached file contains a list of vendors in the Liminal Springs mall\ldots find the vendor that makes the least money relative to the rent it pays. Then, tell me what is listed in the ``type'' column for that vendor.''\\[4pt]
\textbf{Single-hop:} ``Read the Excel file at \texttt{076c8171\ldots.xlsx} and return its contents.''\\
\textbf{Expected substring:} ``Rainforest Bistro.''
\end{quote}

A task is correct iff the expected substring appears (case-insensitive) in the called tool's output. The benchmark contains 109 tasks, 13 to 18 per category.

\subsection{GAIA results}

Table~\ref{tab:gaia-results} reports retriever-only vs.\ retriever+reranker accuracy. The reranker improves overall answer-found rate by $+17.4$\,pp and reduces top-1 synthetic picks from $52\%$ to $3\%$.

\begin{table}[h]
\centering
\small
\begin{tabular}{lrrrrrrrrr}
\toprule
 & Overall & arxiv & excel & fetch & fs & pdf & search & wiki & synth@1 \\
\midrule
Retriever        & $47.7$ & $36$ & $50$ & $22$ & $85$ & $93$ & $33$ & $31$ & $52\%$ \\
+Reranker        & $\mathbf{65.1}$ & $57$ & $0$ & $72$ & $69$ & $71$ & $94$ & $88$ & $3\%$ \\
\bottomrule
\end{tabular}
\caption{GAIA accuracy (\% answer-found), per category and overall. \texttt{synth@1}: fraction of top-1 picks that are synthetic variants.}
\label{tab:gaia-results}
\end{table}

\subsection{Experiment 1: corruption detection}
\label{sec:exp1}

\textit{Goal.} Test whether the reranker learns to downrank our four synthetic failure modes relative to the retriever.

\textit{Setup.} For each GAIA query, score every in-category endpoint with both the retriever and the reranker. We report precision@1 (top in-category scorer is a real tool) and AUC ($P(\text{real} > \text{synth})$ under the same query) by corruption type.

\begin{table}[h]
\centering
\small
\begin{tabular}{lrr}
\toprule
Metric & Retriever & Reranker \\
\midrule
Precision@1 (top in-category is real)   & $0.11$ & $\mathbf{1.00}$ \\
AUC vs \texttt{truncated}                & $0.05$ & $0.39$ \\
AUC vs \texttt{hollow}                   & $0.03$ & $0.51$ \\
AUC vs \texttt{flaky}                    & $0.02$ & $0.41$ \\
AUC vs \texttt{stale}                    & $0.13$ & $0.48$ \\
\bottomrule
\end{tabular}
\caption{Corruption detection. Reranker achieves perfect precision@1 and lifts AUC against every corruption type, though it keeps synthetics reachable below rank 1.}
\label{tab:exp1}
\end{table}

\subsection{Experiment 2: cold start and quality drift}
\label{sec:exp2}

\textit{Goal.} Measure adaptation speed when a tool's quality changes. Three scenarios: a new tool introduced mid-stream, a good tool that degrades, a bad tool that recovers.

\textit{Setup.} A replay simulator of 1{,}500 steps per scenario with a designated canary tool. Quality transitions at step 500. Mean canary rank on a fixed probe set is logged every 50 steps.

\begin{table}[h]
\centering
\small
\begin{tabular}{llrr}
\toprule
Scenario & Transition & Rank ($T{+}250$) & Rank ($T{+}1000$) \\
\midrule
2a: cold start    & introduced at step 500, $95\%$ SR             & $48.0$ & $20.4$ \\
2b: good$\to$bad  & from rank 1.0, flipped to deterministic fail  & $10.2$ & $3.6$ \\
2c: bad$\to$good  & from rank 71.6, flipped to $95\%$ SR          & $70.4$ & $45.6$ \\
\bottomrule
\end{tabular}
\caption{Adaptation trajectories. Downranking (2b) is fast; uptake (2a) is slow; recovery (2c) does not occur without exploration, since an avoided canary is not re-observed.}
\label{tab:exp2}
\end{table}

\subsection{Experiment 3: model-specific tool preferences}
\label{sec:exp3}

\textit{Goal.} Test whether different LLMs require different tools on the same task, and whether the $\operatorname{emb}_{\text{model}}$ channel captures this variation.

\textit{Setup.} For four GAIA probes, force each of the six LLMs onto each of several candidate tools and record whether the string-match oracle finds the expected answer. Each cell is an independent execution.

\begin{table}[h]
\centering
\small
\begin{tabular}{lcccccc}
\toprule
 & gpt-5.4 & grok & gemini & gemma & qwen & deepseek \\
\midrule
\multicolumn{7}{l}{\textbf{arxiv-05}: get arXiv paper 1409.0473} \\
\texttt{arxiv-latex-mcp/get\_paper\_abstract} & pass & pass & pass & pass & pass & pass \\
\texttt{pdf-reader-mcp/read\_pdf}             & \textbf{fail} & pass & \textbf{fail} & \textbf{fail} & pass & pass \\
\midrule
\multicolumn{7}{l}{\textbf{pdf-03}: read PDF contents} \\
\texttt{fabriqa-pdf-reader/read-pdf} & pass & pass & pass & pass & pass & pass \\
\texttt{pdf-reader-mcp/read\_pdf}    & pass & pass & \textbf{fail} & pass & pass & pass \\
\midrule
\multicolumn{7}{l}{\textbf{search-05}: find the year the Berlin Wall fell} \\
\texttt{tavily-mcp/tavily\_search}   & pass & pass & \textbf{fail} & pass & pass & pass \\
\texttt{serper-mcp/google\_search}   & pass & pass & pass & pass & pass & pass \\
\midrule
\multicolumn{7}{l}{\textbf{wikipedia-06}: retrieve Wikipedia article on Carl Nebel} \\
\texttt{wikipedia-mcp-rudra/get\_article}      & pass & pass & pass & pass & pass & pass \\
\texttt{wikipedia-mcp-rudra/search\_wikipedia} & \textbf{fail} & pass & pass & pass & \textbf{fail} & pass \\
\bottomrule
\end{tabular}
\caption{Forced tool $\times$ model grid. In every probe, at least one tool is solvable by some LLMs and unsolvable by others, typically due to how each model fills optional arguments. The reranker's $\operatorname{emb}_{\text{model}}$ reflects these preferences: on the fetch query we measure, a synthetic tool moves from rank 45 (gpt-5.4-nano) to rank 19 (deepseek) based only on the calling-LLM embedding.}
\label{tab:exp3-grid}
\end{table}
